% ============================================================
% Round 10: REINFORCE++ / RLOO
% ============================================================

\subsection{Round 10：REINFORCE++ / RLOO}

\subsubsection{本轮问题}

\begin{conceptbox}
\textbf{学习目标：} 理解 critic-free RL 不只有 GRPO 一条路。REINFORCE++ / RLOO 代表另一类思路：回到基础 policy gradient，用 baseline、KL、clip 和 advantage normalization 获得稳定训练。\\
\textbf{主线位置：}
\[
\text{PPO}
\to
\text{GRPO}
\to
\text{DAPO / Dr.GRPO / GSPO}
\to
\text{REINFORCE++ / RLOO}
\]
\end{conceptbox}

\begin{conceptbox}
\textbf{Highlight：REINFORCE++}\\
\textbf{论文：} \emph{REINFORCE++: Stabilizing Critic-Free Policy Optimization with Global Advantage Normalization}\\
\textbf{arXiv：} \href{https://arxiv.org/abs/2501.03262}{2501.03262}，submitted 2025-01-04，last revised 2025-11-10\\
\textbf{早期标题：} \emph{REINFORCE++: A Simple and Efficient Approach for Aligning Large Language Models}\\
\textbf{Code：} \href{https://github.com/OpenRLHF/OpenRLHF}{OpenRLHF/OpenRLHF}\\
\textbf{关键词：} critic-free RL，global advantage normalization，PPO-clip，token-level KL，REINFORCE++ with baseline\\
\textbf{核心定位：} 不训练 critic，也不依赖 prompt-level local normalization；用全局 advantage normalization 稳定 REINFORCE-style LLM RL。
\end{conceptbox}

\begin{conceptbox}
\textbf{Highlight：RLOO}\\
\textbf{论文：} \emph{Back to Basics: Revisiting REINFORCE Style Optimization for Learning from Human Feedback in LLMs}\\
\textbf{arXiv：} \href{https://arxiv.org/abs/2402.14740}{2402.14740}，submitted 2024-02-22，last revised 2024-02-26\\
\textbf{关键词：} REINFORCE Leave-One-Out，critic-free RLHF，group baseline，online RL at lower cost\\
\textbf{核心定位：} 证明很多 PPO 组件在 LLM RLHF 里未必必要，用更简单的 REINFORCE-style 方法也能有效做在线偏好优化。
\end{conceptbox}

\subsubsection{为什么要回到 REINFORCE}

PPO 的问题是需要 critic：

\[
A_t = R_t - V_\psi(s_t)
\]

GRPO 的做法是去掉 critic，用同一 prompt 的 group reward 做相对 advantage：

\[
\hat{A}_i^\text{GRPO}
=
\frac{r_i-\mu_\text{group}}{\sigma_\text{group}+\epsilon}
\]

但这带来我们前面讲过的问题：

\begin{itemize}[leftmargin=1.5em]
  \item prompt-level / local normalization 可能有偏；
  \item group 太小，advantage 估计不稳；
  \item 全对 / 全错 group 没有信号；
  \item std scaling 可能改变不同 prompt 的训练权重。
\end{itemize}

REINFORCE++ 的问题意识是：

\[
\text{能不能保留 critic-free 的简单性，但不用局部 group std normalization？}
\]

\begin{summarybox}
\textbf{一句话：} REINFORCE++ = 回到 REINFORCE policy gradient，但吸收 PPO 的稳定化技巧，并用 global advantage normalization 替代局部 group normalization。
\end{summarybox}

\subsubsection{REINFORCE 基础公式}

最基础的 REINFORCE policy gradient 是：

\[
\nabla_\theta J(\theta)
=
\mathbb{E}_{y\sim\pi_\theta(\cdot\mid x)}
\left[
R(x,y)
\nabla_\theta
\log \pi_\theta(y\mid x)
\right]
\]

因为：

\[
\log \pi_\theta(y\mid x)
=
\sum_{t=1}^{T}
\log \pi_\theta(y_t\mid x,y_{<t})
\]

所以梯度可以理解成：

\[
R(x,y)
\sum_{t=1}^{T}
\nabla_\theta
\log \pi_\theta(y_t\mid x,y_{<t})
\]

如果 reward 高，就提高这整段回答 token 路径的概率；如果 reward 低，就降低。

为了降低方差，通常加 baseline：

\[
\nabla_\theta J(\theta)
=
\mathbb{E}
\left[
(R-b)
\nabla_\theta
\log \pi_\theta(y\mid x)
\right]
\]

这里：

\[
A = R-b
\]

就是 advantage。

\begin{intubox}
\textbf{直觉：} REINFORCE 不问 critic 这个状态应该拿几分。它只看这条完整回答拿了多少 reward，再用 baseline 判断它比基准好还是差。
\end{intubox}

\subsubsection{RLOO：Leave-One-Out Baseline}

RLOO = REINFORCE Leave-One-Out。

它也会对同一个 prompt 采样多个回答：

\[
y_1,y_2,\ldots,y_G
\]

并得到 reward：

\[
r_1,r_2,\ldots,r_G
\]

对第 $i$ 个回答，baseline 不用 critic，而是用同组其他回答的平均值：

\[
b_i^\text{RLOO}
=
\frac{1}{G-1}
\sum_{j\ne i}
r_j
\]

于是：

\[
A_i^\text{RLOO}
=
r_i-b_i^\text{RLOO}
\]

\begin{summarybox}
\textbf{RLOO 的直觉：} 判断 $y_i$ 好不好时，不拿它自己参与 baseline，而是问：它比同题其他回答平均水平高还是低？
\end{summarybox}

这和 GRPO 很像，都是 group-based critic-free baseline。但区别是：

\[
\text{GRPO 通常用 group mean/std normalization}
\]

\[
\text{RLOO 用 leave-one-out mean baseline}
\]

\subsubsection{REINFORCE++：全局 Advantage Normalization}

REINFORCE++ 的核心强调是 Global Advantage Normalization。

GRPO 的 local normalization 是在每个 prompt 的小 group 内做：

\[
\hat{A}_i^\text{local}
=
\frac{r_i-\mu_\text{group}}{\sigma_\text{group}+\epsilon}
\]

REINFORCE++ 更强调在 global batch 上做归一化。先得到一批样本的 advantage：

\[
A_1,A_2,\ldots,A_N
\]

然后用整个 batch 的均值和标准差归一化：

\[
\hat{A}_i^\text{global}
=
\frac{
A_i-\mu_\text{batch}
}{
\sigma_\text{batch}+\epsilon
}
\]

其中：

\[
\mu_\text{batch}
=
\operatorname{mean}(A_1,\ldots,A_N)
\]

\[
\sigma_\text{batch}
=
\operatorname{std}(A_1,\ldots,A_N)
\]

\begin{intubox}
\textbf{直觉：} GRPO 是“每道题内部统一标尺”；REINFORCE++ 是“整个 batch 用一个统一标尺”。后者希望减少小 group 局部归一化带来的偏置和不稳定。
\end{intubox}

\subsubsection{REINFORCE++ with Baseline}

REINFORCE++ 也有 group-sampling 版本，常被称为 REINFORCE++ with baseline。

对复杂 reasoning 任务，可以先对同一个 prompt 采样多个回答，用 group baseline 得到初始 advantage：

\[
A_i
=
r_i-\mu_\text{group}
\]

然后不在 prompt 内部除以 group std，而是把所有样本放到 global batch 里统一归一化：

\[
\hat{A}_i
=
\frac{
A_i-\mu_\text{batch}
}{
\sigma_\text{batch}+\epsilon
}
\]

这样它结合了两件事：

\begin{itemize}[leftmargin=1.5em]
  \item group baseline：保留同题内部比较；
  \item global normalization：避免每个小 group 自己的 std 改写权重。
\end{itemize}

\begin{summarybox}
\textbf{和 Dr.GRPO 的关系：} Dr.GRPO 也批评 group std scaling。REINFORCE++ with baseline 的思路相近：可以减去 group mean，但不要只靠 prompt-level std normalization；把归一化放到更大的 global batch 上。
\end{summarybox}

\subsubsection{REINFORCE++ 吸收了 PPO 的哪些技巧}

REINFORCE++ 不是裸 REINFORCE。它吸收了 PPO/RLHF 里几个常用稳定化组件：

\begin{itemize}[leftmargin=1.5em]
  \item \textbf{PPO-clip：} 限制 new policy 相对 old policy 的更新幅度；
  \item \textbf{token-level KL penalty：} 在 reward 或 loss 中约束 policy 不要偏离 reference 太远；
  \item \textbf{mini-batch updates：} 对采样数据做小批量、多步更新；
  \item \textbf{advantage normalization：} 用 global batch 标准化降低梯度方差；
  \item \textbf{reward clipping / scaling：} 处理异常 reward，提升数值稳定性。
\end{itemize}

所以它可以写成一种 PPO-style critic-free objective：

\[
\mathcal{J}_\text{R++}(\theta)
=
\mathbb{E}
\left[
\sum_t
\min
\left(
\rho_t(\theta)\hat{A}_t,\;
\operatorname{clip}(\rho_t(\theta),1-\epsilon,1+\epsilon)\hat{A}_t
\right)
-
\beta \operatorname{KL}
\right]
\]

区别在于：

\[
\hat{A}
\text{ 不来自 critic，而来自 reward / baseline / global normalization}
\]

\begin{warnbox}
\textbf{关键点：} REINFORCE++ 不是“完全不用 PPO 技巧”。它是不训练 critic，但保留 PPO-clip、KL、advantage normalization 等稳定化手段。
\end{warnbox}

\subsubsection{REINFORCE++ / RLOO / GRPO 对比}

\begin{center}
{\small
\begin{tabular}{p{2.7cm}p{3.3cm}p{3.4cm}p{3.4cm}}
\hline
 & \textbf{GRPO} & \textbf{RLOO} & \textbf{REINFORCE++} \\
\hline
是否需要 critic
&
不需要
&
不需要
&
不需要
\\
采样方式
&
同 prompt group sampling
&
同 prompt group sampling
&
$k\ge 1$；可单样本，也可 group
\\
baseline
&
group mean
&
leave-one-out group mean
&
可无 baseline，也可 group baseline
\\
归一化
&
local group std normalization
&
通常用 group baseline，形式更简单
&
global batch advantage normalization
\\
主要目标
&
去 critic，用组内相对 reward
&
用 leave-one-out 降低方差
&
去 critic，同时降低 local normalization 偏置
\\
\hline
\end{tabular}
}
\end{center}

\subsubsection{和 Dr.GRPO / GSPO 的关系}

这些方法关注的维度不同：

\begin{center}
{\small
\begin{tabular}{p{3cm}p{9.5cm}}
\hline
\textbf{方法} & \textbf{主要修的问题} \\
\hline
Dr.GRPO
&
批评 GRPO 的 response length normalization 和 reward std scaling bias。
\\
GSPO
&
批评 token-level ratio 和 sequence-level reward 的粒度不匹配。
\\
REINFORCE++
&
批评 prompt-level local normalization 不稳 / 有偏，改用 global advantage normalization。
\\
RLOO
&
用 leave-one-out group baseline 简化 critic-free advantage 估计。
\\
\hline
\end{tabular}
}
\end{center}

\begin{intubox}
\textbf{直觉：} Dr.GRPO 问“归一化有没有引入长度和难度偏置”；GSPO 问“ratio 粒度对不对”；REINFORCE++ 问“advantage 归一化应该在小 group 里做，还是在更大的 batch 上做”。
\end{intubox}

\subsubsection{什么时候更像 GRPO，什么时候更像 REINFORCE++}

如果任务是数学 / 代码，且每个 prompt 能采样多个回答并自动验证，那么 GRPO / DAPO 很自然：

\[
\text{同题多样本}
\to
\text{组内比较}
\]

如果任务更一般，或者每个 prompt 不一定采很多样本，REINFORCE++ 的 $k\ge 1$ 设定更灵活：

\[
\text{采样一批 outputs}
\to
\text{全局 advantage normalization}
\]

如果既想保留同题比较，又想避免 group std bias，可以用 REINFORCE++ with baseline：

\[
A_i = r_i-\mu_\text{group}
\quad
\to
\quad
\text{global normalization}
\]

\subsubsection{本轮最小结论}

\begin{summarybox}
\textbf{REINFORCE++ / RLOO 的核心：}\\
它们和 GRPO 一样，都是 critic-free。不同点在 advantage / baseline 的处理：
\[
\text{GRPO：local group mean/std normalization}
\]
\[
\text{RLOO：leave-one-out group baseline}
\]
\[
\text{REINFORCE++：global advantage normalization}
\]
REINFORCE++ 的主张是：不用 critic，也不用小 group 的局部 std normalization；在全局 batch 上归一化 advantage，可以更稳定、更接近无偏。
\end{summarybox}

\subsubsection{本轮检查题}

\begin{enumerate}[leftmargin=1.5em]
  \item REINFORCE 为什么天然不需要 critic？
  \item RLOO 的 leave-one-out baseline 怎么计算？
  \item GRPO 的 local normalization 和 REINFORCE++ 的 global normalization 有什么区别？
  \item REINFORCE++ 保留了 PPO 的哪些稳定化技巧？
  \item REINFORCE++ with baseline 和 Dr.GRPO 在批评 group std scaling 上有什么相似点？
\end{enumerate}
